% ---------- Parte 7: latencia ----------
\section{Latencia}
\portadilla{7}{Latencia: la figura 6}{Cuánto tarda cada operación, con y sin competencia. Qué se sostiene, un error de compilación que casi se cuela y por qué las medianas con contención bailan tanto.}{Séptima parte: latencia. Ya no medimos cuántas operaciones salen por segundo, sino cuánto tarda cada una.

Veremos las medianas frente al paper, un error de compilación que casi nos hace confirmar una frase del paper por la razón equivocada, las colas de latencia y cuánto podemos fiarnos de una mediana cuando hay muchos hilos.}

\begin{frame}{Cómo se mide la latencia}
  \fuente{de Wolff et al., sección 6.2; EXPLICACION.md, Partes 4 y 6.2}
  \begin{columns}[T,onlytextwidth]
    \column{0.5\textwidth}
    \small\textbf{\color{marino}El experimento}\par\vspace{1mm}
    \footnotesize
    \begin{itemize}
      \item Corren \textbf{p} hilos a la vez. El primero se cronometra; los demás hacen \cod{store} sin parar, para crear contención.
      \item Tres operaciones: \cod{load} (leer), \cod{store} (escribir sin liberar) y \cod{store \& reclaim} (escribir y liberar el objeto viejo).
      \item p = 1 (sin contención), 24 (todos los núcleos) y 56 (más hilos que núcleos).
      \item Binarios oficiales: 5 rondas con Docker y 15 sin Docker.
    \end{itemize}
    \column{0.46\textwidth}
    \small\textbf{\color{marino}Cómo leer una caja}\par\vspace{2mm}
    \centering
    \begin{tikzpicture}[font=\tiny]
      \draw[color=cAutores,line width=0.6pt] (0,0) -- (0,0.55); \draw[color=cAutores,line width=0.6pt] (0,1.75) -- (0,2.5);
      \draw[color=cAutores,line width=0.6pt] (-0.15,0) -- (0.15,0); \draw[color=cAutores,line width=0.6pt] (-0.15,2.5) -- (0.15,2.5);
      \draw[color=cAutores,fill=cAutores!22,line width=0.6pt] (-0.35,0.55) rectangle (0.35,1.75);
      \draw[color=tinta,line width=1pt] (-0.35,1.05) -- (0.35,1.05);
      \node[anchor=west] at (0.45,2.5) {p99: solo el 1\,\% tarda más};
      \node[anchor=west] at (0.45,1.75) {p75};
      \node[anchor=west] at (0.45,1.05) {\textbf{mediana}: la operación típica};
      \node[anchor=west] at (0.45,0.55) {p25};
      \node[anchor=west] at (0.45,0) {p1};
    \end{tikzpicture}\par\vspace{2mm}
    \raggedright\footnotesize La mediana es la experiencia típica; el p99, lo malo que puede ponerse.
  \end{columns}
  \note{Cómo se mide. Se lanzan p hilos. Solo el primero se cronometra; los demás escriben sin parar para crear competencia. Se miden tres operaciones: leer, escribir sin liberar, y escribir liberando el objeto viejo.

  Cada caja resume muchísimas repeticiones. La raya gruesa es la mediana, la operación típica. La caja contiene la mitad central. Los bigotes llegan al 1 por ciento más rápido y al 1 por ciento más lento.

  Recordad la diferencia: la mediana dice cómo va normalmente; el percentil 99, cuánto puede llegar a tardar. Para un algoritmo wait-free, lo segundo es lo que importa.}
\end{frame}

\begin{frame}{Latencia con los binarios oficiales, con Docker}
  \fuente{datos/ultra9/docker/latencia.csv; mediana de 5 rondas de cada percentil}
  \grafica[60mm]{fig6_docker.pdf}
  \note{Esta es nuestra figura 6. Arriba, sin competencia, en escala normal: todos están entre 14 y 44 nanosegundos, y Williams es el más lento en las tres operaciones.

  Abajo, con 24 hilos, en escala logarítmica, donde cada marca vale diez veces más. Autores, Anderson y Folly están juntos, alrededor de 450 nanosegundos en la lectura. Williams y Libstdc++, unas siete u ocho veces por encima.

  Veamos ahora cómo se compara con el paper.}
\end{frame}

\begin{frame}{Medianas: réplica frente al paper}
  \fuente{datos/ultra9/docker/latencia.csv; valores del paper leídos de las etiquetas de su figura 6}
  \footnotesize
  \begin{tabularx}{\linewidth}{@{}p{34mm}CCCCC@{}}
    \toprule
    \textbf{Mediana, ns} & \textbf{Autores} & \textbf{Anderson} & \textbf{Folly} & \textbf{Williams} & \textbf{Libstdc++}\\
    \midrule
    \cod{load}, p = 1 & \textbf{14} & 18 & 15 & 34 & 20\\
    \quad paper & \textbf{23} & 28 & 24 & — & 29\\
    \cod{load}, p = 56 & \textbf{1\,502} & 1\,540 & 2\,977 & 23\,738 & 61\,687\\
    \quad paper & \textbf{1\,140} & 2\,580 & 11\,600 & — & 99\,700\\
    \cod{store}, p = 1 & 22 & 18 & \textbf{17} & 32 & 20\\
    \quad paper & 33 & \textbf{28} & 37 & — & 36\\
    \cod{store}, p = 56 & 1\,516 & \textbf{1\,498} & 4\,813 & 4\,105 & 52\,832\\
    \quad paper & \textbf{1\,470} & 2\,840 & 2\,620 & — & 111\,000\\
    \cod{store \& reclaim}, p = 1 & 28 & 29 & \textbf{23} & 44 & 25\\
    \quad paper & 44 & 44 & 49 & — & \textbf{41}\\
    \cod{store \& reclaim}, p = 56 & \textbf{813} & 1\,499 & 6\,525 & 4\,439 & 47\,198\\
    \quad paper & 1\,570 & \textbf{1\,250} & 5\,060 & — & 59\,300\\
    \bottomrule
  \end{tabularx}\par\vspace{1mm}
  \scriptsize En negrita, la mejor de cada fila. El paper no publica la latencia de Williams. Con p = 24, ver la tabla del anexo.
  \note{Las medianas, frente a las del paper. En cada fila, la mejor va en negrita.

  En la lectura, Autores es el mejor en los dos, con y sin contención. Eso se repite.

  En la escritura sin contención, el paper pone a Autores segundo, tras Anderson. Aquí queda cuarto, tras Folly, Anderson y Libstdc++. Son cinco nanosegundos, un puñado de ciclos, pero el orden es otro.

  Los valores absolutos no se comparan: son máquinas distintas. Comparamos órdenes y proporciones. La gráfica siguiente lo hace de forma visual.}
\end{frame}

\begin{frame}{Cada rival respecto a Autores: paper frente a réplica}
  \fuente{datos/ultra9/docker y nativo\_oficial/latencia.csv; figura 6 del paper; p = 56 en los dos}
  \grafica[58mm]{latencia_relativa.pdf}
  \note{Igual que hicimos con el rendimiento: cada rival dividido entre Autores. Aquí, a la derecha de la raya, Autores es más rápido.

  Arriba, sin contención: todos están muy cerca de la raya. Pequeñas diferencias de nanosegundos, salvo Williams, que va el doble de lento.

  Abajo, con 56 hilos: Libstdc++ está muy a la derecha en los dos, entre 30 y 90 veces más lento que Autores. Anderson, en cambio, en lectura y escritura está pegado a la raya en la réplica, mientras en el paper estaba a la derecha: aquí empatan.

  Las rayitas, sin Docker y con 15 rondas, caen en general cerca de los círculos rellenos.}
\end{frame}

\begin{frame}{Lo que se sostiene en latencia, y lo que no}
  \fuente{datos/ultra9/docker y nativo\_oficial/latencia.csv; EXPLICACION.md, Partes 6.2 y 7}
  \footnotesize
  \begin{tabularx}{\linewidth}{@{}p{56mm}Xp{22mm}@{}}
    \toprule
    \textbf{Lo que dice el paper} & \textbf{Lo que sale aquí, con y sin Docker} & \textbf{Veredicto}\\
    \midrule
    7. El \cod{load} de Autores es siempre el de menor latencia & El más bajo en los tres casos. Con contención, por un 8-12\,\% frente a Folly y Anderson; con p = 56, empate con Anderson & \vsi\ \scriptsize por poco\\
    8. El \cod{store} de Autores es el segundo con p = 1 & Es el cuarto de cinco, por 5 ns & \vno\\
    9. El \cod{store} de Autores es el mejor con contención & Empate con Anderson: 507 frente a 483 ns con p = 24; sin Docker, 461 frente a 470 & \vparte\\
    10. \cod{store \& reclaim} es más lento que \cod{store} con p = 1 & Los cinco & \vsi\\
    11. Con contención, algunos van más rápido en \cod{store \& reclaim} & Con Docker, Autores baja de 1\,516 a 813 ns con p = 56. Sin Docker y 15 rondas, solo un 1-4\,\%: dentro del ruido & \vnoconf\\
    12. La latencia de Williams es casi siempre mayor & Sin contención, la mayor en las 3 operaciones. Con contención supera a Libstdc++ en 6 de 6 y a Folly en 3 de 6 (sin Docker, 5 y 2) & \vparte\\
    \bottomrule
  \end{tabularx}
  \note{Seis afirmaciones del paper sobre latencia, una por fila.

  La séptima, la más importante, se sostiene: la lectura de Autores es la más rápida en los tres casos. Pero con contención, por poco: entre un 8 y un 12 por ciento, y con 56 hilos empata con Anderson.

  La octava no: en escritura sin contención, Autores es cuarto, no segundo. La novena queda a medias: con contención empata con Anderson.

  La décima se cumple: liberar cuesta. La undécima, que con contención liberar acelera, la vimos con Docker, pero al repetir con 15 rondas el efecto desaparece en el ruido: no se confirma.

  Y la duodécima, sobre Williams, solo en parte. Para medirla bien tuvimos que descubrir un error, que es la siguiente diapositiva.}
\end{frame}

\begin{frame}{Una trampa que casi se cuela: el candado escondido de Williams}
  \fuente{datos/ultra9/docker/latencia\_williams\_con\_candado.csv y latencia.csv; EXPLICACION.md, Parte 6.2}
  \begin{columns}[T,onlytextwidth]
    \column{0.5\textwidth}
    \grafica[54mm]{williams_candado.pdf}
    \column{0.47\textwidth}
    \footnotesize
    \textbf{\color{marino}Qué pasó}\par
    La primera medida decía que el \cod{load} de Williams con p = 24 tardaba 40,7 µs, 93 veces más que Autores. Parecía confirmar al paper. Era un error nuestro.\par\vspace{2mm}
    \begin{porque}[: \texttt{-mcx16}]
      Williams guarda 16 bytes y necesita la instrucción \cod{cmpxchg16b}. El compilador solo la usa con \cod{-mcx16}. El binario de latencia del artefacto se compila sin ella, y \cod{boost::atomic} recurre en silencio a un candado interno.
    \end{porque}
    \textbf{\color{marino}Cómo se confirmó}\par
    Mirando el binario: llamadas a \cod{lock\_pool} y ni un \cod{cmpxchg16b}. Con \cod{-mcx16}: 0 llamadas y 10 instrucciones.
  \end{columns}
  \note{Esta diapositiva es una lección de método.

  La primera vez que medimos a Williams en latencia, su lectura con 24 hilos tardaba 40 microsegundos, 93 veces más que Autores. Parecía confirmar la frase del paper. Pero era un error de medida, y era nuestro.

  Williams guarda 16 bytes y necesita una instrucción especial, cmpxchg16b. El compilador solo la usa si se le pide con la bandera -mcx16. El programa de rendimiento la lleva; el de latencia, no, porque el paper deja a Williams fuera de esa prueba. Sin la bandera, la biblioteca Boost mete un candado interno sin avisar. Fue como cronometrar a un corredor con una mochila que nadie le vio poner.

  Lo descubrimos comparando con el programa del equipo, y lo confirmamos mirando las instrucciones del binario.

  Queda una duda: si los autores midieron a Williams así, su frase vendría en parte del candado. No podemos saberlo, porque no publican esos datos.}
\end{frame}

\begin{frame}{Las colas: lo que la mediana no cuenta}
  \fuente{datos/ultra9/docker/latencia.csv; mediana de 5 rondas de cada percentil}
  \grafica[34mm]{colas.pdf}
  \footnotesize
  \begin{columns}[T,onlytextwidth]
    \column{0.48\textwidth}
    \begin{porque}[: la caja de pendientes de Anderson]
      Anderson vacía su caja cada 2\,000 retiradas y recorre las mesas de todos. 1 de cada 2\,000 es el 0,05\,\%: justo lo que mira el p99,95, que salta de 32 a casi 8\,000 ns. No se midió directamente.
    \end{porque}
    \column{0.48\textwidth}
    \begin{porque}[: lock-free no es wait-free]
      Con 23 escritores, el \cod{load} de Williams necesita ganar dos bucles de CAS. El sistema avanza, pero el lector se queda atrás: 1 lectura de cada 2\,000 tarda 10 µs; en Autores, 1,6 µs.
    \end{porque}
  \end{columns}
  \note{Las colas cuentan dos historias que la mediana esconde. Las tres barras de cada algoritmo son la mediana, el percentil 99 y el 99,95, que es lo que solo supera una operación de cada dos mil.

  A la izquierda, sin contención, liberando: todos van finos hasta el 99,95, donde Anderson salta a casi 8 microsegundos. ¿Por qué justo ahí? Anderson vacía su caja de pendientes cada 2 000 retiradas, y uno de cada 2 000 es exactamente el 0,05 por ciento. La coincidencia es demasiado exacta para ser casualidad, aunque no lo medimos directamente.

  A la derecha, con 24 hilos, leyendo: Williams y Libstdc++ tienen colas de 10 y 65 microsegundos. Es lo que significa de verdad lock-free: alguien avanza siempre, pero un hilo concreto puede quedarse atrás. Autores, wait-free, se queda en 1,6.}
\end{frame}

\begin{frame}{¿Cuánto fiarse de una mediana con contención? Poco}
  \fuente{datos/ultra9/docker y nativo\_oficial/latencia.csv; ronda más lenta ÷ más rápida de cada punto}
  \grafica[44mm]{variabilidad.pdf}
  \footnotesize
  \begin{tabularx}{\linewidth}{@{}p{32mm}CCC@{}}
    \toprule
    \textbf{Ronda más lenta ÷ más rápida, mediana} & \textbf{Docker, 5 rondas} & \textbf{Sin Docker, 15 rondas} & \textbf{\cod{bench.cpp}, 5 rondas}\\
    \midrule
    p = 1 & 1,00× & 1,00× & 1,11×\\
    p = 24 & 3,40× (peor, 13,6×) & 5,75× (peor, 36,9×) & 1,84× (peor, 10,4×)\\
    \bottomrule
  \end{tabularx}
  \note{¿Cuánto podemos fiarnos de una mediana? Esta gráfica compara, para cada medida, la ronda más lenta con la más rápida. Uno significa que todas las rondas coinciden.

  Sin contención, todo está en uno: el cronómetro marca siempre lo mismo. Con 24 hilos, la misma medida cambia típicamente entre 3 y 6 veces de una ronda a otra, y en el peor caso, 37 veces.

  ¿Por qué con 15 rondas sale peor que con 5? Porque con más rondas hay más ocasiones de pillar una muy buena y una muy mala, como al tirar un dado más veces. Más rondas no hacen cada ronda más estable; hacen la mediana menos dependiente de la mala suerte.

  Consecuencia: con contención, una diferencia del 20 por ciento entre dos algoritmos no significa nada. Solo valen las enormes, como la de Libstdc++.}
\end{frame}

\begin{frame}{Con p = 24, las rondas caen en dos grupos}
  \fuente{datos/ultra9/nativo\_oficial/latencia.csv, 15 rondas; EXPLICACION.md, Parte 6.2}
  \begin{columns}[T,onlytextwidth]
    \column{0.5\textwidth}
    \grafica[54mm]{bimodal.pdf}
    \column{0.47\textwidth}
    \footnotesize
    \begin{itemize}
      \item Unas rondas quedan por debajo de 520 ns y otras por encima de 1\,000 ns, sin nada en medio.
      \item En el grupo rápido: Autores entre 399 y 445 ns, Folly entre 420 y 517 y Anderson entre 451 y 511.
      \item Autores cae menos veces en el grupo lento: 2 de 15, frente a 6 de Folly y 7 de Anderson. Con 15 rondas no basta para llamarlo propiedad del algoritmo.
    \end{itemize}
    \begin{idea}\footnotesize Con los hilos fijados (\cod{todos24}) el grupo lento no apareció: 419-427 ns en Autores. Indicio de que depende de dónde coloca Linux al hilo medido, en un núcleo P o E.\end{idea}
  \end{columns}
  \note{Mirando las 15 rondas una por una aparece algo curioso: no se reparten de forma continua, sino en dos grupos. Unas por debajo de 520 nanosegundos y otras por encima de mil, sin nada en medio.

  En el grupo rápido, Autores, Folly y Anderson están muy juntos. Autores cae menos veces en el grupo lento, 2 de 15, pero con tan pocas rondas no nos atrevemos a decir que sea una propiedad del algoritmo.

  La pista: cuando fijamos los hilos a núcleos concretos, el grupo lento no aparece. Así que probablemente depende de dónde coloca Linux al hilo que se cronometra, en un núcleo rápido o en uno eficiente. Es un indicio, no una prueba.}
\end{frame}
